\documentclass[12pt,a4paper]{article}

% Packages
\usepackage[T1]{fontenc}
\usepackage[utf8]{inputenc}   % If you're on modern LaTeX, you can remove this.
\usepackage{lmodern}
\usepackage{geometry}
\usepackage{graphicx}
\usepackage{amsmath, amssymb}
\usepackage{caption}
\usepackage{setspace}

% Bibliography (BibLaTeX)
\usepackage[backend=biber,style=numeric,url=true]{biblatex}
\addbibresource{references.bib}

% Hyperlinks (load late)
\usepackage{hyperref}

% Page geometry
\geometry{margin=1in}

% Line spacing
\doublespacing
\setlength{\parindent}{0pt}
\setlength{\parskip}{10pt}

% Title info
\title{Final Exam Paper }
\author{Emma Feege}
\date{\today}

\begin{document}

% Title page
\maketitle

% Abstract
% \begin{abstract}
% A brief summary of your report, highlighting objectives, methods, results, and conclusions.
% \end{abstract}

% Table of contents
\tableofcontents
\newpage

\section{Architecture}

\subsection{Choice of Architectural Style}

The proposed system follows a layered, cloud-native, event-driven microservices architecture. The reason for this is that in the task it was stated that the business will initially be small and cost-constrained but expected to grow very rapidly. Medium-term, this makes scalability, robustness and the ability to independently deploy the required services highly important quality attributes of the system to be developed, which are core strengths of microservices. Microservice architectures also natively support the separate management and parallel development of the individual services, that can hence be developed with the best-suiting technologies and in a more agile way \cite{microsoft_learn_microservices}.

A drawback of this decision in the specific case of the task would be that microservices introduce a complexity overhead compared to more traditional, monolithic approaches. According to the Microsoft Learn information on microservice architectures, topics such as service discovery, data consistency, transaction management and interservice communication pose challenges during the development \cite{microsoft_learn_microservices}. Another aspect to be taken into account is the fact, that transforming to a completely different architecture style is also a big challenge the company. In fact, as Conway's law states, very often, the communication structure of a company is reflected in the architecture/design of the software that the company develops \cite{conways_law}. If this holds true for the enterprise of the task, in the long run, a successful, lasting transition to cloud-native software systems might need more than a well-designed cloud-native architecture for the online business, like organizational and process changes.

\subsection{Proposal for the Solution Architecture}

The proposed system consists of a layered architecture. The individual layers together with the software components they contain will be presented in detail in the following subsections.

\subsubsection{User Experience Layer}

The user experience layer acts as the outermost system layer and is responsible for the direct communication and interaction with the customers of the online business. The frontend application for the catalog and ordering website as well as the blog is hosted using AWS Amplify, which provides managed hosting for web applications. It features deployment, scaling, and content delivery via a Git-based workflow \cite{aws_amplify_docs}.

Furthermore, the user experience layer has a content delivery network (CDN) that caches and delivers static content for the catalog and ordering website and the blog, like HTML pages, CSS files, JavaScript code and images. The reason for including a CDN is the expected latency reduction, which is done by serving the customers from geographically closer locations. This is expected to increase the customer satisfaction. Also, CDNs offer security features like a reduction of the impact of DDoS attacks through load distribution \cite{microsoft_cdns}.

The CDN will be integrated using Amazon CloudFront together with an S3 bucket for the static content and will be automatically provisioned and configured by Amplify. CloudFront offers the needed standard functionality like caching and has different flat-rate plans, that cap the usage costs and would allow a scale-up, as soon as the online business starts to grow \cite{aws_cloudfront_overview}. The S3 bucket acts as the origin location for the static content, from which CloudFront will retrieve it in case it is not cached.

\subsubsection{API Layer}

The communication to the system's backend will be via RESTful API calls, which go through the API layer of the system. The API layer has the purpose of decoupling the user experience layer from the business logic layer, that represents the backend. It consists of an API Gateway, which in particular controls the forwarding of requests to the responsible microservices, manages the traffic and offers monitoring capabilities and access control \cite{aws_apigateway_overview}.

For this, the fully managed Amazon API Gateway is proposed, which has a tiered pricing model, so that over time and when the online business scales up, the price per million individual requests drops  \cite{aws_apigateway_overview}. The authentication and authorization of logged in users (e.g. admins or blog contributors) can be integrated by using Amazon Cognito, which is connected to the frontend via AWS Amplify \cite{aws_integrate_cognito}. The security tokens that are issued by Cognito (authentication) can be validated by the Amazon API Gateway (authorization), before the requests are forwarded further to the backend services \cite{aws_apigateway_cognito}.

\subsubsection{Business Logic Layer}

The business logic layer contains the required core logic of the online business. It consists of the catalog and ordering service, the blog service, the payment service, the customer relationship service and the analytics service.

The microservices are designed as stateless and follow a serverless model. Individually, they are implemented as serverless functions using AWS Lambda. Lambda is preferred over Amazon Fargate, as none of the required services are expected to have long-run running processes or high memory demands \cite{aws_lambda_or_fargate}. AWS Lambda automatically manages the underlying compute infrastructure and scales the services depending on the incoming request volume \cite{aws_lambda}. This makes it well suited for an online business that will likely experience fluctuating workloads, for example due to daily usage patterns, discount promotions, or seasonal demand like the year-end business. In addition, the pay-per-use pricing model of Lambda reduces costs during phases of low activity (e.g. during night-time). The reduced operational overhead will allow the enterprise's development teams to focus on the business functionality instead of infrastructure management for the services that are developed in-house \cite{aws_lambda}.

Services that are provided by third parties are integrated following the same architectural principles, like  modularity, information hiding and loose coupling. Each third party functionality is accessed via an integration microservice, which encapsulates the communication with the external provider. By isolating the external dependencies behind service-internal interfaces, the rest of the system remains independent of the specific third party implementation. This also allows a later replacement of the chosen third-party provider with either another provider or an in-house implementation. A realistic example for this would be the payment service, which could use existing online payment systems such as PayPal internally. The payment service would in this case be responsible for handling API calls to PayPal, authentication with it, data transformations between internal and external formats and error handling.

The event-driven communication between the microservices is implemented using Amazon EventBridge, which enables asynchronous communication between them. For example, when a new order is created, an order event is published to EventBridge, which can trigger downstream services like the payment service, analytics service, or customer relationship service independently.

\subsubsection{Data Layer}

The data layer consists of multiple microservice-specific databases. In order to improve the security and prevent direct public access to sensitive business data, the data layer is deployed inside a virtual private cloud (VPC). Databases are placed in private subnets and are therefore not directly reachable from the internet, but can only be accessed by authorized backend services \cite{aws_vpc}.

For persistent data like orders, the catalog and ordering service uses the managed Amazon RDS with a relational database like MySQL. A relational database is chosen, as order data requires ACID compliance and support for complex queries, for example when retrieving order histories or generating reports \cite{microsoft_rds}. Additionally, hosting the database inside the VPC ensures that access can be restricted using security groups, allowing only the responsible backend services to make connections.

Each microservice is connected to its own dedicated data storage. The catalog and ordering service uses Amazon DynamoDB as a session and state database to store temporary information (e.g. shopping cart contents, checkout progress, and user session data). DynamoDB is fully managed and optimized for high-performance key-value access and automatic scaling, which makes it well suited for session data that frequently changes. Although DynamoDB itself is provided as a managed regional service, access from the backend services can be routed through VPC endpoints so that communication remains within the private AWS network instead of traversing the public internet \cite{aws_vpc}\cite{aws_vpc_dynamodb}.

The blog service and customer relationship service also use Amazon RDS for storing structured and relational data, including blog posts, potentially user accounts, customer requests, and support messages.

\subsubsection{Interface to the existing Fulfillment and Inventory System}

For the existing fulfillment and inventory software system, which is hosted on-premise in the company's internal infrastructure, it is assumed that its main responsibilities are the tracking of products, management of stock levels, order processing, and coordination of shipment and fulfillment. According to the task description, the system follows a traditional three-tier architecture. As it was stated in the task that the company is a "traditional enterprise", it can be assumed that the core logic of the fulfillment and inventory system was not designed according to cloud-native principles, but might rather follow a monolithic design. As a result, migrating and modifying this software might demand a high development effort. For this reason, it is recommended not to migrate the system immediately, but instead to integrate it using an anti-corruption layer. This layer provides an interface between the new cloud-native services and the legacy system and prevents tight coupling between them.

Once the online business grows, migration to the cloud can be considered using the Amazon 6R migration strategies. One option would be to rehost ("lift and shift") the software to the cloud by moving it either into a virtual machine or a container, without changing much. Since the software can stay as is, this would be the least complex, fastest and cheapest option, but might introduce inefficiencies and technical debts into the cloud environment \cite{amazon_6rs}.

Another option would be to replatform ("lift, tinker and shift"), where components such as the database are migrated to managed cloud services while the application logic would remain mostly unchanged. Replatforming would be a more complex and expensive approach, but on the other side would offer an enhanced efficiency due to the partial use of cloud-native functionalities\cite{amazon_6rs}.

In the long term, refactoring or rearchitecting the system into a cloud-native solution may provide the greatest scalability and maintainability benefits, but this would involve significantly higher development effort and cost, so it is not recommended for the initial setup of the online business \cite{amazon_6rs}.

\section{Security, Availability and Performance Aspects}

using the amazon well-architected framework as a reference

\section{Conclusion}

ok

\printbibliography

\end{document}
